In the Final Degree Model, The algoritm has, at all times, the knowledge of the what the degree of all offline edges is in the final graph. While this allows the algorithm to know how many opportunities a vertex will have to be matched, we will sho that this doesn't increase the asymptotic competitiveness of any deterministic algorithm beyond $0.5$
\begin{theorem}
	In the Final Degree model no deterministic algorithm has a competitiveness higher than $0.5$
\end{theorem}
\begin{proof}
	Consider the following graph $G$ which has $n=2k+1$ offline vertices with final degree $k+1$\\
	$k+1$ vertices arrive, all connected to all unmatched offline vertices.\\
	Notice that at this point $k+1$ vertices are matched and the remaining $k$ have degree $k+1$ in the current graph, meaning that they cannot be matched. All that remains is to fulfill the degree promise. This is easy as the degrees of the matched vertices are $1,2,\dots,k,k+1$ meaning that all that is required is for $k$ vertices connected to all offline vertices with unfulfilled degree promise to arrive.\\
	The size of the matching found is $k+1$ while the graph has a perfect matching of size $n=2k+1$ which is achived when the $k$ vertices matched by the graph are instead matched to the vertices which fulfilled their degree promise while the remaining $k+1$ vertices are matched among themselves since they form a biclique $K_{k+1,k+1}$
\end{proof}
\begin{figure}[H]
	\centering
	\scalebox{.8}{\begin{tikzpicture}
	\begin{pgfonlayer}{nodelayer}
		\node [style=default] (0) at (-3, 13) {};
		\node [style=default] (1) at (-3, 12) {};
		\node [style=default] (2) at (-3, 11) {};
		\node [style=default] (3) at (-3, 10) {};
		\node [style=default] (4) at (-3, 9) {};
		\node [style=default] (5) at (-3, 8) {};
		\node [style=default] (6) at (-3, 7) {};
		\node [style=deleted] (7) at (2, 13) {};
		\node [style=deleted] (8) at (2, 12) {};
		\node [style=deleted] (9) at (2, 11) {};
		\node [style=deleted] (10) at (2, 10) {};
		\node [style=deleted] (11) at (2, 9) {};
		\node [style=deleted] (12) at (2, 8) {};
		\node [style=deleted] (13) at (2, 7) {};
		\node [style=none] (14) at (-4, 13) {4};
		\node [style=none] (15) at (-4, 12) {4};
		\node [style=none] (16) at (-4, 11) {4};
		\node [style=none] (17) at (-4, 10) {4};
		\node [style=none] (18) at (-4, 9) {4};
		\node [style=none] (19) at (-4, 8) {4};
		\node [style=none] (20) at (-4, 7) {4};
		\node [style=none] (21) at (2.75, 10) {1};
		\node [style=none] (22) at (2.75, 9) {2};
		\node [style=none] (23) at (2.75, 8) {3};
		\node [style=none] (24) at (2.75, 7) {4};
		\node [style=none] (25) at (2.75, 11) {5};
		\node [style=none] (26) at (2.75, 12) {6};
		\node [style=none] (27) at (2.75, 13) {7};
		\node [style=none] (28) at (-3.25, 6) {final degree};
		\node [style=none] (29) at (2, 6) {arrival order};
	\end{pgfonlayer}
	\begin{pgfonlayer}{edgelayer}
		\draw (6) to (13);
		\draw (5) to (12);
		\draw (4) to (11);
		\draw (3) to (10);
		\draw (2) to (9);
		\draw (1) to (8);
		\draw (0) to (7);
		\draw (0) to (8);
		\draw (0) to (9);
		\draw (1) to (9);
		\draw (1) to (10);
		\draw (2) to (10);
		\draw (2) to (11);
		\draw (3) to (11);
		\draw (3) to (12);
		\draw (4) to (12);
		\draw (4) to (13);
		\draw (4) to (10);
		\draw (5) to (10);
		\draw (5) to (11);
		\draw (5) to (13);
		\draw (6) to (12);
		\draw (6) to (11);
		\draw (6) to (10);
		\draw (0) to (10);
		\draw (1) to (11);
		\draw (2) to (12);
		\draw (3) to (13);
	\end{pgfonlayer}
\end{tikzpicture}
}
\end{figure}